% ---
% file: docs/latex/Fisica/unidades/unidades_fundamentales.tex
% description: Catálogo y especificación de las 7 unidades fundamentales base del SI (redefinición 2019).
% type: doc/manual
% version: 1.0.0
% date: 2026-08-24
% covers:
%   - src/data/units.py
% relations:
%   - docs/fisica/unidades/unidades_principales.md
%   - docs/fisica/unidades/unidades_secundarias.md
%   - docs/ARCHITECTURE.md
% keywords:
%   - unidades-fundamentales
%   - sistema-internacional
%   - unidades-base
%   - constantes-definitorias
% ---

\documentclass[11pt,a4paper]{article}
\usepackage[utf8]{inputenc}
\usepackage[T1]{fontenc}
\usepackage[spanish,es-nodecimaldot,es-noquoting]{babel}
\usepackage{amsmath,amssymb,amsfonts}
\usepackage{esint}
\usepackage{array,tabularx}

% Configuración de márgenes y espaciado estándar
\setlength{\textwidth}{16.5cm}
\setlength{\textheight}{24cm}
\setlength{\oddsidemargin}{-0.25cm}
\setlength{\evensidemargin}{-0.25cm}
\setlength{\topmargin}{-1cm}
\setlength{\parindent}{0pt}
\setlength{\parskip}{5pt}
\setlength{\emergencystretch}{3em}

\begin{document}

\section*{1. Unidades Fundamentales del SI}
\addcontentsline{toc}{section}{1. Unidades Fundamentales del SI}

Este archivo debe leerse primero en la app: define las 7 unidades base del SI (redefinición 2019), ancladas a constantes con valor exacto.

\scriptsize
\begin{tabularx}{\linewidth}{>{\raggedright\arraybackslash}X>{\raggedright\arraybackslash}X>{\raggedright\arraybackslash}X>{\raggedright\arraybackslash}X>{\raggedright\arraybackslash}X>{\raggedright\arraybackslash}X}
\hline
Magnitud fundamental & Unidad base & Símbolo & Constante definitoria asociada & Símbolo de constante & Valor exacto estandarizado (sin incertidumbre) \\
\hline
\textbf{Tiempo} & Segundo & $\mathrm{s}$ & Frecuencia de transición hiperfina del estado fundamental del átomo de $^{133}\mathrm{Cs}$ & $\Delta\nu_{\mathrm{Cs}}$ & $9\,192\,631\,770\ \mathrm{s}^{-1}\ (\mathrm{Hz})$ \\
\hline
\textbf{Longitud} & Metro & $\mathrm{m}$ & Velocidad de la luz en el vacío & $c$ & $299\,792\,458\ \mathrm{m s}^{-1}$ \\
\hline
\textbf{Masa} & Kilogramo & $\mathrm{kg}$ & Constante de Planck & $h$ & $6.626\,070\,15 \times 10^{-34}\ \mathrm{kg m}^2\mathrm{ s}^{-1}\ (\mathrm{J s})$ \\
\hline
\textbf{Corriente eléctrica} & Amperio & $\mathrm{A}$ & Carga elemental & $e$ & $1.602\,176\,634 \times 10^{-19}\ \mathrm{A s}\ (\mathrm{C})$ \\
\hline
\textbf{Temperatura termodinámica} & Kelvin & $\mathrm{K}$ & Constante de Boltzmann & $k_{\mathrm{B}}$ & $1.380\,649 \times 10^{-23}\ \mathrm{kg m}^2\mathrm{ s}^{-2}\mathrm{ K}^{-1}\ (\mathrm{J K}^{-1})$ \\
\hline
\textbf{Cantidad de sustancia} & Mol & $\mathrm{mol}$ & Constante de Avogadro & $N_{\mathrm{A}}$ & $6.022\,140\,76 \times 10^{23}\ \mathrm{mol}^{-1}$ \\
\hline
\textbf{Intensidad luminosa} & Candela & $\mathrm{cd}$ & Eficacia luminosa de radiación monocromática de frecuencia $540 \times 10^{12}\ \mathrm{Hz}$ & $K_{\mathrm{cd}}$ & $683\ \mathrm{cd sr kg}^{-1}\mathrm{ m}^{-2}\mathrm{ s}^3\ (\mathrm{lm W}^{-1})$ \\
\hline
\end{tabularx}
\normalsize


\end{document}
